\documentclass[10pt,a4paper]{amsart}
\usepackage[margin=19mm]{geometry}
\usepackage{amsmath,amssymb,mathtools}
\usepackage[scr=rsfs,cal=euler]{mathalfa}
\usepackage{xfrac}
\usepackage[unicode,hidelinks,bookmarks=false]{hyperref}
\newcommand{\ie}{i.e.\ }
\newcommand{\cf}{cf.\ }
\newcommand{\resp}{respectively\ }
\newcommand{\Subsection}{\subsection}
\providecommand{\nobreakdash}{\nobreak\hbox{-}}
\setlength{\emergencystretch}{2em}
\multlinegap0pt
\usepackage{bm}
\usepackage{bbm}
\usepackage{dsfont}
\usepackage{xspace}
\usepackage{cancel}
\usepackage{mathtools}
\usepackage{amsthm}
\makeatletter
\@ifundefined{Theorem}{\newtheorem{Theorem}{Theorem}}{}
\makeatother
\title[$B'$-orbits on flag varieties and symmetry breaking]{$B'$-orbits on flag varieties and symmetry breaking}
\author{Valentin Massicot}
\date{Source: arXiv:2604.25243v1 (CC BY 4.0)}
\begin{document}
\maketitle

\noindent\emph{Source and licence.} $B'$-orbits on flag varieties and symmetry breaking. Version arXiv:2604.25243v1 (CC BY 4.0), distributed under the Creative Commons Attribution 4.0 International licence, as recorded in the arXiv licence metadata for this exact version and shown on the abstract page https://arxiv.org/abs/2604.25243v1. Full rights notice: licenses/arxiv-2604.25243-CC-BY-4.0.txt.\par\smallskip

\noindent\textit{Editorial scope.} Counted unit: the Theorem whose source label is \texttt{contre exemple} in the article (source0.tex lines 2373--2375, source label \texttt{contre exemple}), delivered together with the article's own complete original proof of it (source0.tex lines 2376--2508, one \texttt{proof} environment of 133 lines). No mathematics is written, repaired or re-derived: statement and proof are the article's own text line for line. Required context retained uncounted: none. Every definition, display and label of those lines is retained unchanged. Omitted from this package and not redistributed: 1--2372, 2509--2859. Nothing inside a retained excerpt is omitted or altered: every retained body is delivered exactly as the article's lines are written, so no omission receipt of any kind accompanies this package. Citations: the retained excerpts carry no citation command at all, so no bibliography entry of the article is transcribed and no reference is redistributed. Reference adaptation: none was needed and none was made; every reference inside the delivered unit resolves inside it, so no unresolved reference or citation is produced. Printed slips kept literal: no printed word, symbol, hypothesis or number of the counted statement or of its proof was altered. Layout and class: the article's own class is \texttt{sigma} with packages including ifpdf, mathtools, amsfonts, amsthm, amssymb, dsfont, enumitem, cleveref, stmaryrd, mathrsfs, makecell, breqn, longtable, tikz; the wrapper is the lane's pinned \texttt{amsart} editorial wrapper at 10pt on A4, which restates the theorem-like environments the retained text needs and adds the article's own macro definitions that the retained text needs (none), with extra wrapper packages (bm, bbm, dsfont, xspace, cancel, mathtools). The delivered numbering is the unit's own. Assets: no retained span contains a figure environment, a graphics-inclusion command, a PDF-page inclusion, a table or a TikZ picture; no image file is copied into this package, the package declares no asset dependency and no image file is redistributed with it. Licence: version arXiv:2604.25243v1 of this article is distributed under the Creative Commons Attribution 4.0 International licence, as recorded in the arXiv licence metadata for this exact version and shown on the abstract page https://arxiv.org/abs/2604.25243v1. A full rights notice is delivered at licenses/arxiv-2604.25243-CC-BY-4.0.txt. Characters: every retained excerpt is pure ASCII, so the delivered engine, XeLaTeX with its default Latin Modern text font, reports no missing character.
\begin{Theorem}{\label{contre exemple}}
	If $n \geqslant 5$, the map $\Psi : B'\backslash G/P_{\underline m} \to \displaystyle \prod_{\substack{L \in \mathcal Q_{B'} \\ H \in \mathcal Q_{P_{\underline m}}}} L\backslash G/H$ is not injective.
\end{Theorem}

\begin{proof}
	Let us first treat the case of $\underline n = (3,2)$ and $\underline m = (1,2,2)$. Consider the flags
	\begin{equation*}
		D_1 = \left[\begin{tabular}{c|cc}
	0 & 1 & 0 \\
	0 & 0 & 1 \\
	1 & 0 & 0 \\
	0 & 0 & 1 \\
	1 & 1 & 0 
\end{tabular}\right], \hspace{3mm} D_2 = \left[\begin{tabular}{c|cc}
	0 & 1 & 0 \\
	0 & 0 & 1 \\
	1 & 0 & 0 \\
	0 & 1 & 1 \\
	1 & 1 & 0 
\end{tabular}\right].
	\end{equation*}
Note that
\begin{equation*}
	\Phi_{\underline m \to (1,4)}(D_1) = \left[\begin{tabular}{c}
	0\\
	0\\
	1\\
	0\\
	1
\end{tabular}\right], \hspace{5mm} \Phi_{\underline m \to (1,4)}(D_2) = \left[\begin{tabular}{c}
	0\\
	0\\
	1\\
	0\\
	1
\end{tabular}\right],
\end{equation*}
\begin{equation*}
	\Phi_{\underline m \to (3,2)}(D_1) = \left[\begin{tabular}{ccc}
	-1 & 1 & 0 \\
	0 & 0 & 1 \\
	1 & 0 & 0 \\
	0 & 0 & 1 \\
	0 & 1 & 0 
\end{tabular}\right] \sim \left[\begin{tabular}{ccc}
	0 & 1 & 0 \\
	0 & 0 & 1 \\
	1 & 0 & 0 \\
	0 & 0 & 1 \\
	0 & 1 & 0 
\end{tabular}\right],
\end{equation*}
\begin{equation*}
	\Phi_{\underline m \to (3,2)}(D_2) = \left[\begin{tabular}{ccc}
	-1 & 1 & 0 \\
	0 & 0 & 1 \\
	1 & 0 & 0 \\
	-1 & 1 & 1 \\
	0 & 1 & 0 
\end{tabular}\right] \sim \left[\begin{tabular}{ccc}
	0 & 1 & 0 \\
	0 & 0 & 1 \\
	1 & 0 & 0 \\
	1 & 0 & 1 \\
	0 & 1 & 0 
\end{tabular}\right] = \left[\begin{tabular}{ccc}
	0 & 1 & 0 \\
	-1 & 0 & 1 \\
	1 & 0 & 0 \\
	0 & 0 & 1 \\
	0 & 1 & 0 
\end{tabular}\right] \sim \left[\begin{tabular}{ccc}
	0 & 1 & 0 \\
	0 & 0 & 1 \\
	1 & 0 & 0 \\
	0 & 0 & 1 \\
	0 & 1 & 0 
\end{tabular}\right]
\end{equation*}
so that $D_1$ and $D_2$ have the same $B'$-invariant ranks.
On the other hand, if we suppose that $D_1$ and $D_2$ are in the same $B'$-orbit then there exists 
\begin{equation*}
	b' = \begin{pmatrix}
	\alpha & \beta & \gamma & 0 & 0 \\
	0 & \delta & \varepsilon & 0 & 0 \\
	0 & 0 & \zeta & 0 & 0 \\
	0 & 0 & 0 & \eta & \theta \\
	0 & 0 & 0 & 0 & \iota
\end{pmatrix} \in B', \hspace{5mm} p = \begin{pmatrix}
	\kappa & \lambda & \mu \\
	0 & \nu & \xi \\
	0 & \rho & \sigma \\
\end{pmatrix} \in P_{\underline m},
\end{equation*}
such that
\begin{equation*}
	b' \begin{pmatrix}
	0 & 1 & 0 \\
	0 & 0 & 1 \\
	1 & 0 & 0 \\
	1 & 0 & 1 \\
	0 & 1 & 0
\end{pmatrix}p = \begin{pmatrix}
	0 & 1 & 0 \\
	0 & 0 & 1 \\
	1 & 0 & 0 \\
	1 & 1 & 1 \\
	0 & 1 & 0
\end{pmatrix},
\end{equation*}
namely
\begin{equation*}
	\begin{pmatrix}
	\ast & \ast & \ast \\
	\kappa \varepsilon & \lambda \varepsilon + \rho \delta & \ast \\
	\ast & \ast & \ast \\
	\kappa \theta & \lambda \theta + \nu \theta + \rho \eta & \ast \\
	\ast & \ast & \ast
\end{pmatrix} = \begin{pmatrix}
	0 & 1 & 0 \\
	0 & 0 & 1 \\
	1 & 0 & 0 \\
	1 & 1 & 1 \\
	0 & 1 & 0
\end{pmatrix}.
\end{equation*}
In particular, one has
\begin{equation*}
	\begin{cases}
	\kappa \theta = 0 \\
	(\lambda+\nu)\theta + \rho \eta = 1 \\
	\kappa \varepsilon = 0 \\
	\lambda \varepsilon + \rho \delta = 0
\end{cases}.
\end{equation*}
Since $\kappa \neq 0$, this implies that $\theta = \varepsilon = 0$ so that $\rho \delta = 0$. In turn, $\delta \neq 0$ so that $\rho = 0$ what is impossible in view of the second equation. Thus, we conclude that $D_1$ and $D_2$ belong to different $B'$-orbits. The theorem then follows for general $\underline n$ and $\underline m$ from observing that $\max(n_1,n_2) \geqslant 3$ whenever $n \geqslant 5$.
\end{proof}

\end{document}
